\documentclass[10pt,a4paper]{book}
\usepackage[utf8]{inputenc}
\usepackage[T1]{fontenc}
\usepackage{lmodern}
\usepackage[a4paper, margin=1.8cm, top=2.2cm, bottom=2.2cm, headheight=16pt, footskip=14pt]{geometry}
\usepackage{xcolor}
\usepackage{graphicx}
\usepackage{tikz}
\usepackage{amsmath,amssymb}
\usepackage{pifont}
\usepackage{fancyhdr}
\usepackage{titlesec}
\usepackage{microtype}
\usepackage{longtable}
\usepackage{booktabs}
\usepackage{array}
\usepackage{tcolorbox}
\usepackage{enumitem}
\usepackage{hyperref}

% Tight options list for questions
\newlist{qoptions}{enumerate}{1}
\setlist[qoptions]{label=\textbf{\Alph*.}, leftmargin=2em, itemsep=1.5pt, topsep=2.5pt, parsep=0pt}

% Color Definitions
\definecolor{hblue}{HTML}{1E3A8A}      % Navy Blue
\definecolor{hred}{HTML}{DC2626}       % Huawei Red accent
\definecolor{singlecol}{HTML}{0D9488}   % Teal for Single Choice
\definecolor{multicol}{HTML}{6D28D9}    % Purple for Multi Choice
\definecolor{tfcol}{HTML}{D97706}       % Amber for True/False
\definecolor{ansgreen}{HTML}{047857}    % Forest Green
\definecolor{ansbg}{HTML}{F0FDF4}       % Light Green Background
\definecolor{qbg}{HTML}{F8FAFC}         % Light Slate Background
\definecolor{qborder}{HTML}{CBD5E1}     % Border Slate
\definecolor{darkslate}{HTML}{1E293B}   % Slate 800

% Hyperref Setup
\hypersetup{
    colorlinks=true,
    linkcolor=hblue,
    citecolor=hblue,
    urlcolor=hblue,
    pdfauthor={Huawei ICT Academy / HCIA-Security Preparation},
    pdftitle={HCIA-Security Comprehensive Exam Question Bank},
    pdfsubject={HCIA-Security V3.0 Certification Exam Preparation},
    pdfkeywords={HCIA, Security, Huawei, Firewall, NAT, VRRP, IPS, PKI, IPsec, VPN}
}

% Header & Footer Configuration
\pagestyle{fancy}
\fancyhf{}
\fancyhead[L]{\small\color{darkslate!70}\textbf{HCIA-Security Exam Bank}}
\fancyhead[R]{\small\color{darkslate!70}\nouppercase{\leftmark}}
\fancyfoot[L]{\small\color{darkslate!50}Huawei Certified ICT Associate - Security}
\fancyfoot[R]{\small\color{darkslate!80}\textbf{Page \thepage}}
\renewcommand{\headrulewidth}{0.6pt}
\renewcommand{\footrulewidth}{0.4pt}

% Section Titles Styling
\titleformat{\part}[display]
  {\normalfont\huge\bfseries\centering\color{hblue}}
  {\Large\MakeUppercase{\partname} \thepart}
  {1ex}
  {\titlerule\vspace{1ex}\Huge}
  [\vspace{1ex}\titlerule]

\titleformat{\section}
  {\normalfont\Large\bfseries\color{hblue}}
  {\thesection}{1em}{}
  [\vspace{0.2ex}{\color{hblue}\titlerule[1.2pt]}]

\titleformat{\subsection}
  {\normalfont\large\bfseries\color{darkslate}}
  {\thesubsection}{0.8em}{}

% Custom Question Box
\newtcolorbox{qcard}[2]{
    colback=qbg,
    colframe=#2,
    arc=1.8mm,
    boxrule=0.7pt,
    leftrule=3.8pt,
    left=7pt,
    right=7pt,
    top=5pt,
    bottom=5pt,
    title={\textbf{Question #1}},
    coltitle=white,
    colbacktitle=#2,
    before=\par\vspace{0.6em}\noindent,
    after=\par\vspace{0.1em}
}

% Custom Explanation Box
\newtcolorbox{explainbox}[1]{
    colback=ansbg,
    colframe=ansgreen,
    arc=1.5mm,
    boxrule=0.5pt,
    leftrule=3.2pt,
    left=7pt,
    right=7pt,
    top=4pt,
    bottom=4pt,
    title={\small\textbf{Correct Answer: #1}},
    coltitle=white,
    colbacktitle=ansgreen,
    before=\par\vspace{0.1em}\noindent,
    after=\par\vspace{0.6em}
}


\begin{document}
\chapter{Single-Answer Questions}
\section{Chapter 1: Network Security Concepts and Specifications}

\begin{qcard}{1}{singlecol}
\textcolor{darkslate!75}{\small\textbf{Topic:} CIA - Confidentiality \hfill \textbf{Type:} Single Choice}\vspace{0.4em}\par
Which security attribute ensures that sensitive information is not disclosed to unauthorized entities or eavesdropped during transmission?
\begin{qoptions}
\item Integrity
\item Confidentiality
\item Availability
\item Non-repudiation
\end{qoptions}
\end{qcard}
\nopagebreak
\begin{explainbox}{B}
\textbf{Concept \& Rationale:} Confidentiality ensures that information can be accessed and viewed only by authorized personnel. Cryptographic encryption is the primary mechanism to enforce confidentiality across untrusted networks.
\end{explainbox}

\begin{qcard}{2}{singlecol}
\textcolor{darkslate!75}{\small\textbf{Topic:} CIA - Integrity \hfill \textbf{Type:} Single Choice}\vspace{0.4em}\par
Which security property guarantees that transmitted data has not been altered, deleted, or forged in transit?
\begin{qoptions}
\item Availability
\item Integrity
\item Confidentiality
\item Authenticity
\end{qoptions}
\end{qcard}
\nopagebreak
\begin{explainbox}{B}
\textbf{Concept \& Rationale:} Integrity guarantees that data has not been tampered with or modified during transit or storage. Cryptographic hash functions (such as SHA-256) and message authentication codes are primary mechanisms to verify integrity.
\end{explainbox}

\end{document}
